\section{Conclusion}
\label{sec:conclusion}

By reformulating the combinatorial constraints of the Discretizable Distance Geometry Problem (DDGP) in the language of linear algebra over $\mathbb{F}_2$, this work establishes a bridge between local geometric reflections and global counting complexity. The core insight is that the complexity of solution counting under pruning is fundamentally governed by the algebraic interaction of labeled violations. When physical violations are systematically indexed by their corresponding mirror cliques, the space of feasible branch shifts emerges naturally as the projection $\mathcal{K}_F = M(\ker V)$. Whenever a viable reference solution exists and mirror-separation holds, this linear structure determines the total solution count by assigning an independent factor of two to each branch decision that escapes the algebraic reach of the pruning constraints.

This perspective provides a unified explanation for the tractability of classical DMDGP and unpruned DDGP instances, where graph-induced symmetries naturally guarantee full branch viability. In the presence of arbitrary pruning edges, the labeled violation matrix $V$ acts as a precise algebraic filter, identifying exactly which combinations of reflections preserve distance constraints. Crucially, our resulting rank formula cleanly decouples the purely combinatorial structure of the discretization graph from the measure-zero geometric coincidences that characterize degenerate physical embeddings.

Our algebraic framework opens a promising avenue for the design of highly efficient distance geometry algorithms. By shifting the focus from combinatorial search over exponential branch trees to linear operations over $\mathbb{F}_2$, this formulation provides the mathematical foundation for solver paradigms that may reduce branch enumeration. In particular, the algebraic filters developed here can be directly integrated into solvers to detect infeasible branches at early stages, bypassing expensive coordinate computations.

Specifically, future work will focus on exploiting the structured sparsity of the mask matrix $M$ and the labeled violation matrix $V$ to construct and solve the associated $\mathbb{F}_2$ systems efficiently for large-scale discretization networks. Such advancements hold the potential to significantly accelerate distance geometry solvers, demonstrating the power of finite-field linear algebra to drastically reduce the computational complexity of spatial reconstruction problems.